\section{Relation to existing thermodynamic inference}
\label{sec:pa-prior-art}

The central practical question is how much dissipation is forced by a specified
observation law. Optimizing over hidden models compatible with observations is
an established approach. Ehrich constructs compatible masked Markov models and
minimizes their entropy production \cite{ciEhrich2021}; the inversion assumes
the correct hidden-state number. Nitzan, Ghosal and Bisker formulate an
optimization over a specified CTMC topology using observed transitions and
waiting-time moments \cite{ciNitzan2023}. Their moment hierarchy progressively
tightens the compatible-model bound. Their numerical experiments with a
smaller hidden topology do not establish a universal reduction to that topology;
the paper explicitly leaves that question open.

This distinction matters for interpreting computations. A feasible physical
model provides an upper bound on the \emph{minimum compatible dissipation}.
It need not bound the dissipation of the actual system from above. Conversely,
a local numerical minimum in a selected architecture need not provide a
valid lower bound over a larger physical model class. The direct correlation
certificate developed here and the solvable circulant optimization serve
different purposes: the former is valid without a specified hidden graph,
whereas the latter resolves an explicitly restricted realization problem.

\subsection{The particular role of a scalar finite-lag defect}

Reversible autocorrelation functions have positive spectral measures, and their
log-convexity follows from Cauchy--Schwarz. Those facts predate this work.
Thermodynamic information in correlation functions is also well established.
Ohga, Ito and Kolchinsky relate cross-correlation asymmetry to cycle affinity
\cite{rfOhga2023}. Shiraishi bounds fluctuation oscillations by entropy
production \cite{rfShiraishi2023}. Van Vu, Vo and Saito extend entropy and
affinity bounds to finite-time cross-correlations \cite{paVu2024}.
In particular, their equation~(13) bounds a squared cross-asymmetry divided by
autocorrelation decay. Their Appendix~D, equations~(112)--(117), uses the
contraction of path relative entropy to the joint endpoint law. Raux, Maier
and Seifert give another arbitrary-lag bound for two-observable
cross-asymmetry \cite{rfRaux2026}; the consulted version is a preprint.

With one common observable, the numerator of each of these two-observable
asymmetry tests vanishes. The finite-lag construction here instead uses
\[
 \Delta_s=-s^2 C(0)+2sC(t)-C(2t).
\]
Writing the stationary transition operator as $P_t=H_t+B_t$, with $H_t$
self-adjoint and $B_t$ skew-adjoint, gives
\[
 \Delta_s=\|B_t f\|_\pi^2-\|(H_t-sI)f\|_\pi^2.
\]
Thus a quantity available from one scalar autocorrelation bounds an
unobserved antisymmetric response from below. The subsequent conversion to
dissipation uses established relative-entropy contraction and bounded-response
inequalities. The contribution to be assessed is this particular finite-lag
combination and its calibrated confidence procedure. It is neither a new
data-processing inequality nor the first use of correlations for thermodynamic
inference.

The sharper nonlinear conversion also has an established statistical ingredient.
With the convention
$J(P,Q)=[D(P\Vert Q)+D(Q\Vert P)]/2$, Sason's Proposition~4 states
\[
 \min_{\operatorname{TV}(P,Q)=z}J(P,Q)=2z\operatorname{atanh}z
 \qquad (0\leq z<1)
\]
\cite{paSason2015}. Applying this inequality to conditional transition laws,
then averaging with the appropriate convexity argument, is distinct from
claiming the divergence inequality itself. General sharp-divergence and
thermodynamic-response results provide further context
\cite{ciGuntuboyina2014,ciFalasco2022}.

Time-symmetric scalar fluctuations already support other entropy estimators.
The variance sum rule and its proofs derive a lower bound using the initial
curvature of displacement variance in overdamped Langevin dynamics, with
calibrated diffusion coefficients \cite{ciVSR2024,ciVSRProofs2024}. Full
force-assisted measurements can yield an equality. The present bounded
observable uses specified nonzero lags and a stationary Markov-semigroup
identity; it does not require estimating a derivative at zero. These are
different access assumptions, and neither method should be described as
universally stronger without a comparison under matched access and errors.

\subsection{Circulant optima and the limits of a spectral interpretation}

The thermodynamic cost of coherent oscillations has an extensive literature.
The affinity-and-topology constraint of Barato and Seifert
\cite{rfBarato2017} and its proof by Ohga and coworkers
\cite{rfOhga2023} must be distinguished from an entropy-production bound
using a relaxation eigenvalue alone. Gu proves
\[
 \sigma\geq \eta\,\frac{\omega^2}{\gamma},
 \qquad
 \eta=\frac{\|v\|_{2,\pi}^2}{\|v\|_\infty^2},
\]
for an oscillatory eigenmode $v$ with eigenvalue $-\gamma+i\omega$
\cite{rfGu2026}. In that paper, Section~V explicitly treats arbitrary
translation-invariant jump processes on a ring, establishes $\eta=1$ for
Fourier modes, and gives the nearest-neighbor ring formulas. The drift-diffusion
limit on a circle attains $\sigma=\omega^2/\gamma$ there. That continuum
envelope and its saturation are established results. The finite-size
minimization over all allowed circulant displacements is a more specific
optimization problem, whose proof is supplied in this work.

Kolchinsky derives a dissipation-coherence inequality from a condition on
phase-current fluctuations and proves it in the weak-noise Gaussian regime,
with some extensions beyond Gaussian fluctuations \cite{paKolchinsky2026}.
This is not an unconditional eigenvalue-only theorem for arbitrary finite
Markov networks. Shiraishi's counterexample retains a finite oscillatory
eigenvalue while making its driven part increasingly rare, with stationary
entropy production tending to zero \cite{rfShiraishi2026}. The consulted
counterexample is an unrefereed preprint. It illustrates why observable
amplitude and stationary population cannot be removed from a universal
inference claim. The scalar $\Delta_s$ preserves those weights.

\subsection{Temporal access, uncertainty and implementation}

Sampling a continuous-time system at a fixed interval can leave its generator
ambiguous. System aliasing and the associated matrix-logarithm branches are
established identification problems \cite{paYue2021}. The circle diffusion
construction makes their thermodynamic consequence explicit: different drifts
can produce the same sampled transition kernel while dissipating very
different amounts. A finite set of irregular lags requires a separate
finite-precision analysis. Approximate aliases from simultaneous Diophantine
approximation concern a fixed finite timing menu and a bounded record length;
they do not automatically establish an impossibility result for continuously
randomized timing, unlimited records or every adaptive experiment. Randomized
clocks under detector deadtime are studied in the recent spectral-access
preprint \cite{rfAznagulov2026}.

The analog acquisition scheme combines standard conditional-variance
identities, bounded-observation concentration and direct use of the measured
record. Its value is an explicit interface between the detector contract,
calibration uncertainty and a physical lower bound. A squared noisy reading
has second moment $B=\mathbb{E}Z^2$, while the latent conditional mean has
$A=\mathbb{E}(\mathbb{E}[Z\mid X])^2\leq B$. They are different quantities.
The sign of the calibration coefficient makes that distinction usable in a
conservative joint inequality, while retaining its potential loss of power.
The implementation and acquisition-matched simulations evaluate this complete
procedure; they do not replace laboratory validation.

The primary sources checked in this audit do not state the precise scalar
finite-lag entropy certificate and joint analog procedure presented here.
That finding supports a focused contribution claim, while leaving historical
priority to broader specialist review. A formal result within the stated
stationary classical models does not establish a new law for arbitrary
quantum, inertial or nonstationary systems.
